% !TEX root = ../main.tex
% ---------------------------------------------------------------------
%  Navigation: Farblegende + "Wenn die Aufgabe lautet ..."
%  Steht auf der Titelseite, vor dem Inhaltsverzeichnis.
% ---------------------------------------------------------------------

\begin{tcolorbox}[zfbox, colback=black!2, colframe=black!45,
  coltitle=white, colbacktitle=black!45, title={Farbcode der Blöcke}]
\textcolor{cB1}{\rule{7pt}{7pt}}~\T{B1} Grundlagen \quad
\textcolor{cB2}{\rule{7pt}{7pt}}~\T{B2} TwinCAT/SPS \quad
\textcolor{cB3}{\rule{7pt}{7pt}}~\T{B3} ST \quad
\textcolor{cB4}{\rule{7pt}{7pt}}~\T{B4} Filter \\
\textcolor{cB5}{\rule{7pt}{7pt}}~\T{B5} Automaten \quad
\textcolor{cB6}{\rule{7pt}{7pt}}~\T{B6} Tasks/EtherCAT \quad
\textcolor{cB7}{\rule{7pt}{7pt}}~\T{B7} PID \quad
\textcolor{cB8}{\rule{7pt}{7pt}}~\T{B8} Safety \\[3pt]
\textcolor{cREF}{\rule{7pt}{7pt}}\,Formel \;
\textcolor{red!65}{\rule{7pt}{7pt}}\,Falle \;
\textcolor{green!40!black}{\rule{7pt}{7pt}}\,Rezept \;
\textcolor{black!55}{\rule{7pt}{7pt}}\,Beispiel \\[3pt]
\T{MEP:} 120\,min, 120\,Punkte (1\,Punkt $\approx$ 1\,Minute), alle
Papierunterlagen erlaubt, Code \emph{auf Papier}. Ja/Nein-Fragen:
falsch $=-1$ \ra bei Unsicherheit leer lassen.
Abschnitt~11 ist ein Stichwortregister.
\end{tcolorbox}

\begin{tcolorbox}[zfbox, colback=black!2, colframe=black!45,
  coltitle=white, colbacktitle=black!45, title={Wenn die Aufgabe lautet \dots}]
\begin{tabularx}{\linewidth}{@{}Xr@{}}
Ja/Nein-Fragebogen (HS24, mit Begründung) & \ref{sec:quiz} \\
Variable in GVL mit Klemme verknüpfen & \ref{rez:gvlio} \\
Sensorspannung \ra physikalische Grösse & \ref{rez:skalierung} \\
Warum reicht eine \R{FUNCTION}? & \ref{sec:pou} \\
$G(s)$ mit $s=\frac{1-z^{-1}}{T}$ diskretisieren & \ref{rez:diskret} \\
Differenzengleichung als FB programmieren & \ref{bsp:fbfilter} \\
Zustandsdiagramm skizzieren, ENUM + CASE & \ref{rez:automat} \\
Förderband-Automat komplett (MEP HS24) & \ref{bsp:foerderband} \\
Ampel mit Timern, Hand/Automatik & \ref{bsp:ampel} \\
Zweite Anlage mit gleicher Logik & \ref{falle:instanz} \\
Timer (\R{TON}/\R{TOF}/\R{TP}) richtig einsetzen & \ref{sec:timer} \\
Flanke / Doppelklick erkennen & \ref{bsp:edge} \\
Statuswort zerlegen, \R{SHR}/\R{SHL}, Hex & \ref{rez:statuswort} \\
Gleitender Mittelwert / Forgetting Factor & \ref{sec:mittelwert} \\
PID diskretisieren (P, I-Trapez, D gefiltert) & \ref{rez:pid} \\
Anti-Reset-Windup (Back-Calculation) & \ref{sec:arw} \\
Tasks: Basiszeit, Priorität, Zeitverlauf & \ref{bsp:tasks} \\
EtherCAT: on-the-fly, in-place, WC & \ref{sec:ethercat} \\
RETAIN vs.\ PERSISTENT & \ref{sec:remanenz} \\
\end{tabularx}
\end{tcolorbox}
